\documentclass[10pt,landscape]{article}
\usepackage[utf8]{inputenc}
\usepackage[landscape,top=0.43in,bottom=0.43in,left=0.35in,right=0.35in]{geometry}
\usepackage{multicol}
\usepackage{amsmath,amssymb,amsthm}
\usepackage{enumitem}
\usepackage{xcolor}
\usepackage{titlesec}

% Compact spacing
\setlength{\parindent}{0pt}
\setlength{\parskip}{0pt}
\setlength{\itemsep}{0pt}
\setlength{\topsep}{0pt}
\setlength{\partopsep}{0pt}

% Compact lists
\setlist{nosep,leftmargin=*,topsep=0pt,partopsep=0pt}

% Compact sections with better contrast
\titlespacing{\section}{0pt}{4pt}{2pt}
\titlespacing{\subsection}{0pt}{3pt}{1pt}
\titleformat{\section}{\normalfont\fontsize{12}{13}\bfseries\color{blue!70!black}\centering}{\thesection}{1em}{}
\titleformat{\subsection}{\normalfont\fontsize{9.5}{10.5}\bfseries}{\thesubsection}{1em}{}

% Colored boxes for key concepts
\newcommand{\bbox}[1]{\colorbox{yellow!20}{{\small\textbf{#1}}}}

\pagestyle{empty}

\begin{document}

\begin{center}
    \Large{\textbf{DS 3000: Foundations of Data Science - Formula Sheet}} \\
    \small{Midterm Exam Reference}
\end{center}

\begin{multicols*}{3}
\setlength{\columnseprule}{0.4pt}

% ============= PAGE 1, COLUMN 1 =============
\section*{Matrices vs Vectors}

\textbf{Vector:} A column of numbers (one-dimensional array)
$$\vec{v} = \begin{bmatrix} v_1 \\ v_2 \\ \vdots \\ v_n \end{bmatrix} \text{ (size: } n \times 1)$$

\textbf{Matrix:} A rectangular array of numbers (two-dimensional)
$$A = \begin{bmatrix} a_{11} & a_{12} & \cdots & a_{1n} \\ a_{21} & a_{22} & \cdots & a_{2n} \\ \vdots & \vdots & \ddots & \vdots \\ a_{m1} & a_{m2} & \cdots & a_{mn} \end{bmatrix}$$

\textbf{Dimensions:} $m \times n$ means $m$ rows by $n$ columns

\textbf{Remember:} Rows first, columns second (RC)

% ============= SECTION 1 =============
\section*{Vector Operations}

\subsection*{Basic Operations}
\textbf{Addition:} $\vec{u} + \vec{v} = \begin{bmatrix} u_1 + v_1 \\ u_2 + v_2 \\ \vdots \\ u_n + v_n \end{bmatrix}$

\textbf{Scalar Multiplication:} $c\vec{v} = \begin{bmatrix} cv_1 \\ cv_2 \\ \vdots \\ cv_n \end{bmatrix}$

\textbf{Element-wise (Hadamard):} $\vec{u} \odot \vec{v} = \begin{bmatrix} u_1v_1 \\ u_2v_2 \\ \vdots \\ u_nv_n \end{bmatrix}$

\subsection*{Dot Product}
$\vec{u} \cdot \vec{v} = u_1v_1 + u_2v_2 + \cdots + u_nv_n = \vec{u}^T\vec{v}$

\textbf{Example:} $\begin{bmatrix} 2 \\ 3 \end{bmatrix} \cdot \begin{bmatrix} 4 \\ -1 \end{bmatrix} = 2(4) + 3(-1) = 8 - 3 = 5$

Result is a \textbf{scalar} (single number)

\subsection*{Vector Magnitude (L2-Norm)}
$\|\vec{v}\| = \sqrt{v_1^2 + v_2^2 + \cdots + v_n^2} = \sqrt{\vec{v} \cdot \vec{v}}$

\columnbreak

% ============= PAGE 1, COLUMN 2 =============
\section*{Matrix Operations}

\subsection*{Matrix Dimensions}
For matrix $A_{m \times n}$:
\begin{itemize}
    \item $m$ = number of rows (first)
    \item $n$ = number of columns (second)
    \item Entry $a_{ij}$ is in row $i$, column $j$
\end{itemize}

\subsection*{Transpose}
$(A^T)_{ij} = A_{ji}$ (swap rows and columns)

For $A_{m \times n}$, result $A^T$ is $n \times m$

\subsection*{Matrix-Vector Multiplication}
$A\vec{x} = \begin{bmatrix} \text{row}_1 \cdot \vec{x} \\ \text{row}_2 \cdot \vec{x} \\ \vdots \\ \text{row}_m \cdot \vec{x} \end{bmatrix}$

For $A_{m \times n}$ and $\vec{x}_{n \times 1}$, result is $m \times 1$

\textbf{Dimension rule:} Inner dimensions must match

\textbf{Example:} $\begin{bmatrix} 1 & 2 \\ 3 & 4 \end{bmatrix} \begin{bmatrix} 5 \\ 6 \end{bmatrix} = \begin{bmatrix} 1(5)+2(6) \\ 3(5)+4(6) \end{bmatrix} = \begin{bmatrix} 17 \\ 39 \end{bmatrix}$

Each entry is dot product of row with $\vec{x}$

\subsection*{Matrix Multiplication}
$(AB)_{ij} = \sum_{k=1}^{n} A_{ik}B_{kj}$

For $A_{m \times n}$ and $B_{n \times p}$, result is $m \times p$

\textbf{Dimension rule:} $(m \times \mathbf{n})({\mathbf{n}} \times p) = (m \times p)$

\textbf{Not commutative:} $AB \neq BA$ (generally)

\textbf{Associative:} $(AB)C = A(BC)$

\subsection*{Element-wise (Hadamard)}
$(A \odot B)_{ij} = A_{ij}B_{ij}$

Requires $A$ and $B$ to have same dimensions

\subsection*{Identity Matrix}
$I_n = \begin{bmatrix} 1 & 0 & \cdots & 0 \\ 0 & 1 & \cdots & 0 \\ \vdots & \vdots & \ddots & \vdots \\ 0 & 0 & \cdots & 1 \end{bmatrix}$

$AI = IA = A$

\columnbreak% ============= PAGE 1, COLUMN 3 =============
% ============= SECTION 4 =============
\section*{Systems of Linear Equations}

\subsection*{Matrix Form}
$A\vec{x} = \vec{b}$

Augmented matrix: $[A \; | \; \vec{b}]$

\subsection*{Row Operations}
\begin{enumerate}
    \item Swap two rows
    \item Multiply a row by nonzero scalar
    \item Add multiple of one row to another
\end{enumerate}

\subsection*{Row Echelon Form (REF)}
\begin{itemize}
    \item All zero rows at bottom
    \item Leading entry of each row is to right of leading entry above
    \item Leading entries are nonzero
\end{itemize}

\subsection*{Reduced Row Echelon Form (RREF)}
REF plus:
\begin{itemize}
    \item Leading entry in each row is 1
    \item Leading 1 is only nonzero entry in its column
\end{itemize}

\subsection*{Solution Types with RREF Examples}

\textbf{Unique solution:} Pivot in every column of $A$
$$\left[\begin{array}{cc|c} 1 & 0 & 3 \\ 0 & 1 & -2 \end{array}\right] \implies x_1 = 3, x_2 = -2$$

\textbf{Infinite solutions:} Free variables (columns without pivots)
$$\left[\begin{array}{ccc|c} 1 & 0 & 2 & 4 \\ 0 & 1 & -1 & 3 \\ 0 & 0 & 0 & 0 \end{array}\right]$$
$x_1 = 4 - 2t, x_2 = 3 + t, x_3 = t$ (free)

\textbf{No solution:} Row $[0 \; 0 \; \cdots \; 0 \; | \; b]$ where $b \neq 0$
$\left[\begin{array}{cc|c} 1 & 0 & 2 \\ 0 & 1 & 5 \\ 0 & 0 & 1 \end{array}\right] \implies 0 = 1 \text{ (inconsistent)}$

\newpage

% ============= PAGE 2, COLUMN 1 =============
\section*{Projections}

\subsection*{Projection onto Vector}
\textbf{Projection of $\vec{b}$ onto $\vec{a}$:}
$$\text{proj}_{\vec{a}}\vec{b} = \frac{\vec{a} \cdot \vec{b}}{\vec{a} \cdot \vec{a}} \vec{a} = \frac{\vec{a} \cdot \vec{b}}{\|\vec{a}\|^2} \vec{a}$$

This gives the point in span($\vec{a}$) closest to $\vec{b}$

\textbf{Scalar component:} $\text{comp}_{\vec{a}}\vec{b} = \frac{\vec{a} \cdot \vec{b}}{\|\vec{a}\|}$

\subsection*{How to Choose Variables}
\textbf{Problem format:} ``Find the point in the span of the $\vec{a}$ vectors closest to the $\vec{b}$ vector''

\textbf{Rule:} Project the vector you want to get close to (b) \textit{onto} the span of the other vector(s) (a)

$$\text{proj}_{\vec{a}}\vec{b} \text{ means: project } \vec{b} \text{ onto } \vec{a}$$

Bottom subscript = the vector defining the subspace

Top = the vector being projected

% ============= SECTION 6 =============
\section*{Linear Independence}

\subsection*{Definition}
Vectors $\vec{v}_1, \vec{v}_2, \ldots, \vec{v}_k$ are \textbf{linearly independent} if:
$$c_1\vec{v}_1 + c_2\vec{v}_2 + \cdots + c_k\vec{v}_k = \vec{0}$$
implies $c_1 = c_2 = \cdots = c_k = 0$

\subsection*{Testing Independence}

\textbf{Method 1: Solve homogeneous system}

Form matrix $A = [\vec{v}_1 \; \vec{v}_2 \; \cdots \; \vec{v}_k]$

Solve $A\vec{x} = \vec{0}$:
\begin{itemize}
    \item \textbf{Independent:} Only trivial solution ($\vec{x} = \vec{0}$)
    \item \textbf{Dependent:} Nontrivial solutions exist
\end{itemize}

\textbf{Method 2: Determinant (for square matrices)}

For $n$ vectors in $\mathbb{R}^n$, form matrix $A = [\vec{v}_1 \; \vec{v}_2 \; \cdots \; \vec{v}_n]$:
\begin{itemize}
    \item $\det(A) \neq 0$ $\implies$ \textbf{independent}
    \item $\det(A) = 0$ $\implies$ \textbf{dependent}
\end{itemize}

\textbf{Quick checks:}
\begin{itemize}
    \item More vectors than dimensions → dependent
    \item Contains zero vector → dependent
    \item Two vectors are scalar multiples → dependent
\end{itemize}

\columnbreak

% ============= PAGE 2, COLUMN 2 =============
\section*{Eigenvalues \& Eigenvectors}

\subsection*{Definition}
$A\vec{v} = \lambda\vec{v}$ where $\vec{v} \neq \vec{0}$

$\lambda$ is eigenvalue, $\vec{v}$ is eigenvector

\subsection*{Characteristic Equation}
$$\det(A - \lambda I) = 0$$

\textbf{For 2×2 matrix:}
$$A = \begin{bmatrix} a & b \\ c & d \end{bmatrix}$$
$$\det(A - \lambda I) = (a-\lambda)(d-\lambda) - bc = 0$$
$$\lambda^2 - (a+d)\lambda + (ad-bc) = 0$$

\textbf{Trace:} $\text{tr}(A) = a + d = \lambda_1 + \lambda_2$

\textbf{Determinant:} $\det(A) = ad - bc = \lambda_1 \lambda_2$

\subsection*{Finding Eigenvectors}
For each eigenvalue $\lambda$:
\begin{enumerate}
    \item Solve $(A - \lambda I)\vec{v} = \vec{0}$
    \item Find null space of $(A - \lambda I)$
\end{enumerate}

% ============= SECTION 8 =============
\section*{Determinants}

\subsection*{2×2 Matrix}
$$\det\begin{bmatrix} a & b \\ c & d \end{bmatrix} = ad - bc$$

\subsection*{3×3 Matrix (Cofactor Expansion)}
$$\det\begin{bmatrix} a & b & c \\ d & e & f \\ g & h & i \end{bmatrix}$$

Expand along first row:
$$= a\det\begin{bmatrix} e & f \\ h & i \end{bmatrix} - b\det\begin{bmatrix} d & f \\ g & i \end{bmatrix} + c\det\begin{bmatrix} d & e \\ g & h \end{bmatrix}$$

$$= a(ei - fh) - b(di - fg) + c(dh - eg)$$

\textbf{Pattern:} Alternate signs ($+, -, +$) along row/column

Can expand along any row or column

\columnbreak

% ============= PAGE 2, COLUMN 3 =============
\section*{Linear Combinations \& Span}

\subsection*{Linear Combination}
$\vec{v} = c_1\vec{v}_1 + c_2\vec{v}_2 + \cdots + c_k\vec{v}_k$

where $c_i$ are scalars

\subsection*{Span}
$\text{span}(\vec{v}_1, \vec{v}_2, \ldots, \vec{v}_k) = \{c_1\vec{v}_1 + c_2\vec{v}_2 + \cdots + c_k\vec{v}_k : c_i \in \mathbb{R}\}$

\textbf{Check if $\vec{b}$ in span:} Solve $A\vec{x} = \vec{b}$ where $A = [\vec{v}_1 \; \vec{v}_2 \; \cdots \; \vec{v}_k]$

$\vec{b}$ is in span if system is consistent (has solution)

\section*{Properties Summary}

\subsection*{Dot Product Properties}
\begin{itemize}
    \item $\vec{u} \cdot \vec{v} = \vec{v} \cdot \vec{u}$ (commutative)
    \item $\vec{u} \cdot (\vec{v} + \vec{w}) = \vec{u} \cdot \vec{v} + \vec{u} \cdot \vec{w}$ (distributive)
    \item $(c\vec{u}) \cdot \vec{v} = c(\vec{u} \cdot \vec{v})$ (scalar mult)
\end{itemize}

\subsection*{Transpose Properties}
\begin{itemize}
    \item $(A^T)^T = A$
    \item $(A + B)^T = A^T + B^T$
    \item $(AB)^T = B^TA^T$ (reverse order)
    \item $(cA)^T = cA^T$
\end{itemize}

\subsection*{Matrix Multiplication Properties}
\begin{itemize}
    \item Associative: $(AB)C = A(BC)$
    \item NOT commutative: $AB \neq BA$ (generally)
    \item Distributive: $A(B+C) = AB + AC$
    \item Identity: $AI = IA = A$
\end{itemize}

\subsection*{Determinant Properties}
\begin{itemize}
    \item $\det(AB) = \det(A)\det(B)$
    \item $\det(A^T) = \det(A)$
    \item $\det(cA) = c^n\det(A)$ for $n \times n$ matrix
    \item $\det(A^{-1}) = \frac{1}{\det(A)}$
    \item If triangular, $\det(A) =$ product of diagonal entries
\end{itemize}

\subsection*{Eigenvalue Properties}
\begin{itemize}
    \item $\det(A) = \prod \lambda_i$ (product of eigenvalues)
    \item $\text{tr}(A) = \sum \lambda_i$ (sum of eigenvalues)
    \item Eigenvectors for distinct $\lambda$ are linearly independent
    \item If $A$ is triangular, eigenvalues are diagonal entries
\end{itemize}

\end{multicols*}

\end{document}